% ============================================================================
% Numerical demonstration section for paw_isdf_thc_prb.tex (AS SHIPPED).
% Inserted AFTER \section{Error analysis} (sec:error), BEFORE sec:scaling.
% Fig.~1 = Result B (ISDF-rank convergence).
% NOTE 2026-06-10: the earlier Result-A ERI-accuracy table (tab:eri-accuracy)
% and its subsection were REMOVED per user direction (implementation detail,
% not a fundamental ISDF statement). The A numbers live in paw_results_A.csv
% as code validation only -- do NOT re-add to the paper.
% ============================================================================
\section{Numerical demonstration}
\label{sec:numerics}

We demonstrate that the PAW-ISDF-THC factorization converges systematically to
the exact electron repulsion integrals as the ISDF auxiliary basis is enlarged,
and that the all-electron augmentation preserves this convergence.  All
calculations use the implementation in the \textsc{CoQu\'i} code, taking
Quantum~ESPRESSO mean-field data through \texttt{pw2coqui}.  The reference
system is LiH in the zinc-blende-derived rocksalt structure on a
\(2\times2\times2\) \(\bk\)-mesh (three irreducible \(\bk\)-points, sixteen
bands), treated with three pseudization families that share the same valence:
norm-conserving (NCPP), ultrasoft (USPP), and projector-augmented-wave (PAW).
This lets us isolate the effect of the augmentation algebra while holding the
smooth ISDF problem fixed.

Equation~\eqref{eq:error-bound} bounds the ERI error in the Coulomb metric by
the interpolation error of the compensated smooth pair densities, controlled by
the number of global interpolation points \(N_\mu\).  Figure~\ref{fig:convergence}
shows the relative Frobenius error of the Hartree and exchange matrix elements
against the direct reference as the ISDF threshold is tightened from
\(3\times10^{-2}\) to \(10^{-6}\), sweeping \(N_\mu\) from roughly \(120\) to
\(430\).  The error falls smoothly and nearly as a power law in \(N_\mu\),
spanning about two and a half orders of magnitude for all three pseudization
families.  Two features are notable.  First, the ultrasoft and PAW curves are
essentially superimposed: the local augmentation rows in the composite
auxiliary index converge at the same rate as the smooth problem and do not
introduce a slower-decaying error channel.  Second, the norm-conserving curve
is shifted to smaller \(N_\mu\)---it reaches a given accuracy with fewer
interpolation points---because its pseudo pair densities require no compensation
charges, whereas the ultrasoft and PAW cases add a modest amount of auxiliary
content to represent the augmentation.  This is precisely the behavior
anticipated by Eq.~\eqref{eq:error-bound}: the augmented factorization inherits
the controlled, tunable convergence of the underlying momentum-resolved ISDF,
with a small and systematically convergent augmentation overhead.

\begin{figure*}[t]
  \includegraphics[width=0.92\textwidth]{paw_convergence_B.pdf}
  \caption{Relative Frobenius error of the tensor-hypercontracted Hartree
  (left) and exchange (right) matrix elements versus the number of ISDF
  interpolation points \(N_\mu\), for LiH on a \(2\times2\times2\) mesh with
  norm-conserving (NCPP), ultrasoft (USPP), and projector-augmented-wave (PAW)
  pseudization.  Tightening the ISDF threshold from \(3\times10^{-2}\) to
  \(10^{-6}\) drives the error down by about two and a half orders of
  magnitude.  PAW converges at the same rate as the ultrasoft case (the two
  curves overlap), while the norm-conserving case reaches a given accuracy at
  smaller \(N_\mu\) because it carries no compensation charges.}
  \label{fig:convergence}
\end{figure*}
